\begin{table}[htbp]
\raggedright
\caption{Profile moderators of attacked effectivity: controlled contrasts, SEM paths, and descriptive susceptibility weights.}
\label{tab:multivariate}
{
\footnotesize
\setlength{\tabcolsep}{4pt}
\renewcommand{\arraystretch}{1.08}
\resizebox{\linewidth}{!}{%
\begin{tabular}{p{0.34\linewidth}rrrrrrrrr}
\toprule
Moderator & Role & Controlled mean b & Controlled p & Ridge mean b & Weight \% & SEM mean b & SEM mean |b| & SEM min p & SEM paths \\
\midrule
Big Five Neuroticism Mean \% & core & 2.013 & 0.214 &  &  & 2.086 & 2.644 & 0.051 & 4.000 \\
Sex Other & core & -7.216 & 0.064 & 0.154 & 2.783 & -5.143 & 11.557 & 0.143 & 4.000 \\
Big Five Conscientiousness Mean \% & core & -1.925 & 0.280 &  &  & -1.888 & 1.888 & 0.349 & 4.000 \\
Sex Female & core & 5.636 & 0.110 & 1.301 & 6.894 & 2.813 & 13.692 & 0.001 & 4.000 \\
Age Years & core & 0.082 & 0.968 &  &  & -0.204 & 2.083 & 0.372 & 4.000 \\
Big Five Openness To Experience Mean \% & core & 0.073 & 0.971 &  &  & 0.105 & 1.876 & 0.387 & 4.000 \\
Big Five Extraversion Mean \% & exploratory & -0.336 & 0.846 &  &  &  &  &  &  \\
Big Five Agreeableness Mean \% & exploratory & -0.308 & 0.852 &  &  &  &  &  &  \\
\bottomrule
\end{tabular}
}
}
\par\smallskip
{\normalsize Note. Controlled coefficients summarize moderator associations with mean attacked effectivity. SEM columns summarize the repeated-outcome path model across the attacked opinion leaves. Weight percentages come from the target-conditional regularized aggregation used to compute the post hoc susceptibility index.}
\end{table}